% !PRINTED 608
\weberpage{608}
%<<CONT
die Wurzeln der Gleichung
\weq{2}{x^p - a = 0.}

Ist diese Gleichung reducibel, so ist
\weq{3}{x^p - a = f_1(x) f_2(x),}
und $f_1$, $f_2$ sind Functionen in $\Omega$ von niedrigerem als dem $p^{\text{ten}}$ Grade.
Ein Theil der Wurzeln (1) wird auf $f_1 = 0$, ein anderer Theil auf $f_2 = 0$ fallen. Ist
also $\mu$ der Grad von $f_1$, so muss für irgend einen Exponenten $\lambda$
\[ \varepsilon^\lambda \alpha^\mu = b \]
eine Grösse in $\Omega$ sein [das von $x$ unabhängige Glied in $f_1(x)$], also, wenn man
in die $p^{\text{te}}$ Potenz erhebt,
\weq{4}{a^\mu = b^p.}

Nun ist $0 < \mu < p$ und daher $\mu$ und $p$ relativ prim, so dass sich zwei ganze Zahlen
$h$ und $k$ aus der Gleichung
\[ \mu h + p k = 1 \]
bestimmen lassen. Es ist dann also nach (4)
\[ a = a^{\mu h} a^{p k} = (b^h a^k)^p; \]
also wäre, gegen die Voraussetzung, $a$ die $p^{\text{te}}$ Potenz einer Grösse in $\Omega$.

Dieser Satz ist, wie man sieht, unabhängig von der Voraussetzung, dass $\Omega$ ein
reeller Körper sei. Nehmen wir ihn aber reell an und ist $p$ nicht gleich 2, so hat die
Gleichung (2) imaginäre Wurzeln, und kann nach 1. eine Normalgleichung $g(t) = 0$ nicht
reduciren. Wenn aber $p = 2$ ist, dann kann $g(t)$ nur dann durch diese Gleichung (2)
reducirt werden, wenn $a$ positiv und der Grad von $g(t)$ eine \gesperrt{gerade} Zahl ist; also:

2.~\gesperrt{Eine Normalgleichung ungeraden Grades kann nicht durch Adjunction eines
reellen Radicals reducirt werden.}

Dieser Fall trifft bei dem casus irreducibilis der cubischen Gleichung zu, wenn man die
Quadratwurzel aus der Discriminante dem Rationalitätsbereich der Co\"efficienten
adjungirt hat. Denn da die Discriminante positiv ist, so bleibt der Rationalitätsbereich
reell, und er bleibt auch reell, wenn man noch so viele reelle Radicale adjungirt. Soll
die Gleichung also durch reelle Radicale lösbar sein, so muss sie bei solchen
Adjunctionen endlich zerfallen, was nach 2. unmöglich ist.
